\section{What the Simulator Looks Like}
\label{sec:landscape}
The simulator was released after the offline phase, and the TA confirmed that local runs may be used to choose models,
policies and their constants, provided the submitted code uses only the training rows before its first call and embeds
no locally measured point, value or weight. We used this in two studies. Both inform the design of
Section~\ref{sec:strategy}; neither feeds data into it.

\textbf{The surrogate cannot see the low region.} On held-out training rows our additive ARD kernel beats generic
21-dimensional ARD kernels (MAE 3.40 against 4.74--5.98), and its 95\,\% intervals cover 94.8\,\% of held-out rows. Its
constant prior mean, however, fails away from the data: the 200 lowest of 6,548 locally measured points ($R\approx105.1$,
about 1.9 training spacings out, on the simplex boundary) are predicted at $124.4\pm8.8$. Trend terms, a log target and
per-building length scales did not change this. Reaching the low region is an exploration problem, not a modelling one.

\textbf{Structure of $R$.} An unrestricted local search (about 22,500 simulator evaluations, exploration only) found
$R=102.10$ and four regularities: (i) $R$ is additive over buildings within a basin (error about 0.03) but couples by up
to 0.85 across basins, and each building has several near-equivalent valleys; (ii) the ramping weight is the sensitive
coordinate and $R$ is terraced, with plateaus 0.02--0.2\,\% of a weight wide; (iii) small weights act as switches:
shares of $10^{-5}$ to $3\times10^{-3}$ on otherwise unused objectives lower $R$ by flipping near-tied hourly decisions,
while shares below $10^{-6}$ equal zero; (iv) the last gains came from systematic one-coordinate scans that random
perturbations missed. Additivity was already visible in the training rows: a building-additive model cross-validates as
well as an unconstrained one ($R^2$ 0.846 vs.\ 0.844).
